\documentclass[11pt]{article}
\usepackage[margin=0.9in]{geometry}
\usepackage{amsmath,amssymb,amsthm}
\usepackage[hidelinks]{hyperref}
\newtheorem{theorem}{Theorem}
\newcommand{\meas}{\operatorname{meas}}
\newcommand{\Li}{\operatorname{Li}_2}
\title{The limiting nonunimodular-root ratio for $P_{m,1}$\\
\large A solution to MathDB 353931}
\author{Principia Math In Progress}
\date{18 September 2026}
\begin{document}
\maketitle
\begin{abstract}
We prove Stankov's conjecture that the limit ratios associated with
$P_{m,1}$ tend to approximately $0.209$. Starting from the published
fixed-$m$ root-counting formula, a cellwise comparison proves convergence
with error at most $2/(3m)$. The limit is
$2\bigl(\Li(1/2)-\Li(-1/2)\bigr)/\pi^2
=0.20885431502703392737\ldots$.
\end{abstract}

\section*{Task 1: the original problem and its full solution}
Write $\Phi_m(x)=\sum_{j=0}^{m-1}x^j$. The family in the source is
\[
 P_{m,1}(x,y)=x^{m-1}\Phi_m(x)+x^{m-1}y+\Phi_m(x)y^2.
\]
For integers $m\geq2$ and $n\geq1$, set
\begin{align*}
 H_{n,m}(x)
 &=\Phi_m(x)\bigl(x^{2n+m-1}+1\bigr)+x^{n+m-1}\\
 &=x^{n+m-1}\left(1+\sum_{j=0}^{m-1}
                 (x^{n+j}+x^{-n-j})\right).
\end{align*}
This is the source's $H_{2n+2m-2}$, of degree $2n+2m-2$, and
$P_{m,1}(x,x^{n+m-1})=x^{m-1}H_{n,m}(x)$.
Let $C(H)$ be the number of roots of $H$ off the unit circle,
counted with multiplicity, divided by $\deg H$, and put
\[
 L_m=\lim_{n\to\infty}C(H_{n,m}).
\]
The fixed-$m$ limit and its following measure representation are
\emph{prior results of Stankov}, not claims of this note:
\begin{equation}\label{measure}
 L_m=\frac1\pi\meas\{t\in[0,\pi]:2|\sin(mt)|<\sin t\}.
\end{equation}
See Theorem 1 and the proof of Theorem 9 in the published paper
\cite{S}, especially page 429. Its envelope is
$\sin(mT/2)/\sin(T/2)$ and its central coefficient is $1$;
substitution $T=2t$ gives \eqref{measure}.
The monomial factor $x^{m-1}$ changes the root count and degree
by a fixed amount for fixed $m$, so it does not change this limit.
Thus $L_m$ is precisely the $LC(P_{m,1})$ of Conjecture 11
on page 431 (Conjecture 2.1 in the arXiv version).
The conjecture asks whether $L_m$ converges to approximately $0.209$.

\begin{theorem}\label{main}
For every integer $m\geq2$,
\begin{equation}\label{error}
 |L_m-L|\leq\frac{2}{3m},
\end{equation}
where
\begin{align}
 L&=\frac{4}{\pi^2}\int_0^{\pi/2}
                    \arcsin\!\left(\frac{\sin t}{2}\right)\,dt
       \label{integral}\\
  &=\frac{2}{\pi^2}\bigl(\Li(1/2)-\Li(-1/2)\bigr)
   =\frac{2}{\pi^2}\sum_{k=0}^{\infty}
                    \frac{4^{-k}}{(2k+1)^2}.
       \label{value}
\end{align}
In particular, the conjectured limit exists and equals
\[
 L=0.2088543150270339273715943563264722403065\ldots.
\]
\end{theorem}

\begin{proof}
Define $h(s)=2\arcsin(s/2)/\pi$ for $0\leq s\leq1$ and
$g(t)=h(\sin t)$ for $0\leq t\leq\pi$.
Partition $[0,\pi]$ into the $m$ intervals
\[
 I_j=[j\pi/m,(j+1)\pi/m],\qquad 0\leq j<m.
\]
For a constant threshold $s\in[0,1]$, the substitution
$u=mt-j\pi$ shows that
\begin{equation}\label{frozen}
 \frac{m}{\pi}\meas\{t\in I_j:2|\sin(mt)|<s\}=h(s).
\end{equation}
Indeed, the set of $u\in[0,\pi]$ with $2\sin u<s$ has length
$2\arcsin(s/2)$, also when $s=0$.

Let $u_j$ and $v_j$ be respectively the minimum and maximum of
$\sin t$ on $I_j$. Monotonicity in the threshold and
\eqref{frozen} give
\[
 h(u_j)\leq A_j:=\frac{m}{\pi}\meas
 \{t\in I_j:2|\sin(mt)|<\sin t\}\leq h(v_j).
\]
The cell average $B_j=(m/\pi)\int_{I_j}g(t)\,dt$ lies in the
same interval. Consequently
\begin{align*}
 \left|L_m-\frac1\pi\int_0^\pi g(t)\,dt\right|
 &\leq\frac1m\sum_{j=0}^{m-1}|A_j-B_j|\\
 &\leq\frac1m\sum_{j=0}^{m-1}
                  \bigl(\max_{I_j}g-\min_{I_j}g\bigr)
 \leq\frac{\operatorname{Var}_{[0,\pi]}g}{m}.
\end{align*}
Here each oscillation is at most the variation on its cell,
and variations add over a partition. The function $g$ increases
from $0$ to $h(1)=1/3$ on $[0,\pi/2]$ and decreases to $0$
on $[\pi/2,\pi]$. Its total variation is therefore $2/3$.
This proves \eqref{error}; reflection about $\pi/2$ gives
\eqref{integral}. It also proves convergence without assuming
independence of $t$ and $mt$ or interchanging the two root-counting
limits.

It remains to evaluate the integral. For $0\leq a\leq1/2$, let
$J(a)=\int_0^{\pi/2}\arcsin(a\sin t)\,dt$.
The derivative of the integrand is uniformly bounded by
$2/\sqrt3$ on this rectangle, so differentiation under the integral
is justified. For $0<a\leq1/2$, substitution $u=\cos t$ gives
\[
 J'(a)=\int_0^1\frac{du}{\sqrt{1-a^2+a^2u^2}}
      =\frac{\operatorname{arctanh}a}{a}
      =\sum_{k=0}^{\infty}\frac{a^{2k}}{2k+1}.
\]
The last series converges uniformly for $0\leq a\leq1/2$.
Since $J(0)=0$, termwise integration yields
\[
 J(a)=\sum_{k=0}^{\infty}\frac{a^{2k+1}}{(2k+1)^2}
     =\frac{\Li(a)-\Li(-a)}2,
 \qquad \Li(z)=\sum_{r=1}^{\infty}\frac{z^r}{r^2}.
\]
Putting $a=1/2$ proves \eqref{value}.
\end{proof}

\paragraph{Endpoints and numerical value.}
All sets above are measured by Lebesgue length, so shared cell
endpoints have no effect. Equality in the threshold also has
measure zero: the analytic functions $2\sin(mt)\pm\sin t$ have
only finitely many zeros on this compact interval and are not
identically zero. The positive series in \eqref{value} has the
explicit tail bound
\[
 0<\sum_{k=N}^{\infty}\frac{4^{-k}}{(2k+1)^2}
 \leq\frac{4^{1-N}}{3(2N+1)^2}.
\]
The accompanying script uses $N=80$ and rational bounds on $\pi$
from Machin's formula to certify the displayed forty decimal
digits. Numerical root-location checks are supplementary only.

\section*{Task 2: the sharp rate of convergence}
Determine the optimal order of $|L_m-L|$, and the leading term
of $L_m-L$. In particular, establish or refute an $O(m^{-2})$
bound and decide whether even and odd $m$ have different leading
coefficients. The estimate in Theorem~\ref{main} does not answer
these questions.

\section*{Task 3: monotonicity and finite-$m$ inequalities}
Determine whether $L_m$ is strictly increasing for every $m\geq2$
and whether $L_m<L$ always holds. Seek explicit finite-$m$ bounds
that improve both the source's estimates and \eqref{error}.

\section*{Task 4: a variable central coefficient}
Replace $x^{n+m-1}$ in $H_{n,m}$ by $\lambda x^{n+m-1}$, with
$\lambda\geq0$. Determine the outer limit as a function of
$\lambda$, including the transition when the threshold
$\lambda\sin t/2$ exceeds $1$. Obtain exact expressions and
regularity statements at the transition.

\section*{Task 5: simultaneous growth of the two parameters}
Find quantitative conditions on $n=n(m)$ ensuring that
$C(H_{n(m),m})\to L$. This requires bounds for the inner
root-counting limit that are uniform in $m$. Investigate the
regime where $n/m$ tends to a finite positive constant as well
as the regime $n/m\to\infty$.

\paragraph{Verification and attribution.}
The fixed-$m$ formula \eqref{measure} is credited to \cite{S};
related general counting and averaging methods occur in
\cite{BM,S25a,S25b}. This note proves the particular outer limit,
its explicit value, and \eqref{error}. The documented bounded
search found no prior full claim for that target. All four
followups are proposed questions, not claimed results. Review
was by the authoring agent only.

\begin{thebibliography}{9}
\raggedright
\small
\setlength{\itemsep}{0pt}
\bibitem{S} D. Stankov, \emph{The number of nonunimodular roots
of a reciprocal polynomial}, Comptes Rendus Math\'ematique
\textbf{361} (2023), 423--435.
\href{https://doi.org/10.5802/crmath.422}{doi:10.5802/crmath.422}.
Earlier version: \href{https://arxiv.org/abs/2203.07710}{arXiv:2203.07710v1} (2022).
\bibitem{BM} D. W. Boyd and M. J. Mossinghoff,
\emph{Small limit points of Mahler's measure}, Experimental
Mathematics \textbf{14}(4) (2005), 403--414, Section 6.
\href{https://doi.org/10.1080/10586458.2005.10128936}{doi:10.1080/10586458.2005.10128936}.
\bibitem{S25a} D. Stankov, \emph{The alternative to Mahler measure
of a multivariate polynomial},
\href{https://arxiv.org/abs/2502.02803}{arXiv:2502.02803v1} (2025).
\bibitem{S25b} D. Stankov, \emph{Approximation of the number of
nonunimodular zeros of a self-reciprocal polynomial},
\href{https://arxiv.org/abs/2509.01015}{arXiv:2509.01015v1} (2025).
\end{thebibliography}
\end{document}
